\section{A simultaneous family of metrics}
\label{sec:geometry}

The positive bridges require a range of self-dual periods.  We construct
that range inside the interval cube, while keeping precisely the genuine
cuts prescribed in Section~\ref{sec:stack}.  Two distinctions matter in this
construction.  First, a vanishing coefficient in a period is an auxiliary
plumbing degeneration; it is not the middle three-sphere cut of an
interval.  Second, the narrow-period region is controlled by a long tube
with spherical fibers.  We retain that tube while replacing the toric
metric near the corner where both coefficients vanish.

Fix a finite stack, its determinant and spin data, and its finitely many
bit choices.  The number of intervals, all exceptional evaluations, and
the topological charge ranges are fixed in every limiting assertion below.
Constants for compact pieces may depend on this finite problem.  The
constant attached only to an old cap will be identified separately in
Section~\ref{sec:estimates}.  A \emph{genuine length} is a length tending to
infinity at a face of the interval cube.  The only such cuts are the
middle spheres $J_i\cong S^3$ and the boundaries of the $(-4)$ disk bundles
$N_i$.  All other lengths used here, including the external surgery seams
$Y_j$, will ultimately be chosen finite.  The $Y$ collars varied below
are the external rational seams joining the ordinary factors: macros,
single intervals, and caps.

In a macro containing positive bridge $j$, orient the adjacent spheres as
\begin{equation}
 X_1=-S_j,\qquad X_2=S_{j+1}.
 \label{eq:geometry:oriented-spheres}
\end{equation}
Write $U$ for the sphere in Proposition~\ref{prop:stack:lattice}.  Thus
$X_i^2=-4$, $U^2=0$, $U X_i=1$, and $X_1X_2=0$.  The regular neighborhoods
needed below are
\begin{equation}
 C_2=\nu(X_1\cup U\cup X_2),\qquad
 C_{1,i}=\nu(U\cup X_{3-i}),\qquad C_0=\nu U,\qquad N_i=\nu X_i.
 \label{eq:geometry:neighborhoods}
\end{equation}
Here neighborhoods are rounded, and their separating hypersurfaces are
specified up to isotopy.  In particular $C_0$ is a neighborhood supporting
a tube, not an additional infinite cut.

\begin{theorem}[The metric family]
\label{thm:geometry:family}
For the fixed finite stack, there are families with the following
properties.  They are defined for all sufficiently small $\epsilon>0$,
all sufficiently small regularizers $d>0$, and arbitrarily long finite
auxiliary rational collars.  The auxiliary bounds may depend on
$\epsilon$ and the fixed problem.  The family parameter is the original
cube of interval parameters, with no added boundary faces.
\begin{enumerate}[label=\textup{(\roman*)}]
\item At its genuine faces the family has independent product length
collars at $J_i$ and $\partial N_i$, compatible at intersections of
faces.  The $J_i$ cross-sections are round.  On a component separated by
$J_i$, the limiting data depend only on the interval parameters retained
on that component.  The $N_i$ metric is complete with positive scalar
curvature and a positive-scalar-curvature product end.  Its link and its
scale on the stretching collar are independent of the tangential cube
parameters.  The link is joined to the round lens metric through
positive-scalar-curvature metrics.  These metric choices are common to
the two bit terms compared at that lens face.
\item On the positive-width toric region of a macro constructed below,
as $d\to0$ the retained main piece has a complete
metric conformal to a toric K\"ahler metric with positive cohomology ray represented by
\begin{equation}
 H=U+aX_1+b_0X_2,\qquad
 0\le a,b_0\le\tfrac14,\qquad a+b_0>0.
 \label{eq:geometry:period}
\end{equation}
It has nonnegative scalar curvature, with strict positivity on an open
set.  A zero coefficient removes the corresponding outer part by a
$C_{1,i}$ separation; a coefficient $1/4$ separates $N_i$.  The class
$H$ is supported rationally on the retained main piece in all these
cases.  Its positive rank is one.
\item Every macro in the small-width replacement region contains an
$S^2\times S^1\times I$ tube, with fiber $U$, of length comparable to
$\epsilon^{-1}$.  On its bulk the metric is, or approaches smoothly,
\begin{equation}
 g_{\mathrm{long}}
 =g_{S^2(1/\sqrt2)}+\tfrac12\,d\xi^2+2\,d\theta_x^2.
 \label{eq:geometry:long-model}
\end{equation}
The same description holds for every toric macro whose normalized
length $L=1/(a+b_0)$ tends to infinity.  All such tube lengths are $O(\epsilon^{-1})$.
The tubes lie in $C_0$ and are disjoint from both middle spheres $J_i$.
\item Along any sequence of these parameters, every additional
unbounded region is one of the following: an exact rational product
collar, an annulus approaching a positive-scalar-curvature product, or
one of the tubes in \textup{(iii)}.  The remaining pieces, after the
specified satellite rescalings, have uniformly bounded geometry on
fixed truncations.  There are only finitely many such pieces for the
fixed stack.  Away from fixed-length tube ends, the error from
\eqref{eq:geometry:long-model}, with any fixed number of spatial
derivatives, tends to zero as the distance from both ends tends to
infinity.  In particular the integrated errors used in
Section~\ref{sec:estimates} are $o(\epsilon^{-1})$, uniformly under
deeper regularization and longer auxiliary rational collars.
For scalar-curvature costs, the middles of the $Y_j$ collars
are omitted, retaining prefixes of length $O(\epsilon^{-1/2})$;
no positivity is required on a $Y_j$ cross-section.
\item For each fixed real cohomology class $\Lambda$ there are closed
comparison forms with its prescribed vertical flux on these tubes,
zero curvature on the long parts of rational collars, and
$o(\epsilon^{-1})$ energy away from the tubes.  More precisely, the
uniform comparison and error assertions are those of
Lemma~\ref{lem:geometry:comparison}.
\end{enumerate}
The construction is smooth at finite parameter values and has the
ordinary product description on every genuine length collar.  Only
spatial derivative estimates are asserted uniformly as lengths diverge;
no bound on derivatives in a compactifying length coordinate is needed.
\end{theorem}

The proof occupies the rest of this section.  The theorem supplies the
geometry for the simultaneous Abelian tests, whose statement and proof
are Theorem~\ref{thm:estimates:clamp}.  Harmonic background connections
and their exact product convention at genuine faces are constructed in
Section~\ref{sec:estimates}; this will also ensure that the complementary
data agree under the lens flip.

\subsection{Embedded plumbings and the period polygon}

\begin{lemma}[Placement of the interfaces]
\label{lem:geometry:plumbing}
The neighborhoods in \eqref{eq:geometry:neighborhoods} can be placed
strictly inside the macro so that $U$ is disjoint from $J_1\cup J_2$,
$C_{1,i}$ is disjoint from $J_i$, $N_{3-i}\subset C_{1,i}$, and
$C_0\subset C_{1,1}\cap C_{1,2}$.  The simultaneous separating
interfaces required below can be chosen disjoint or nested.  The two
interfaces $\partial C_{1,1}$ and $\partial C_{1,2}$ are never required
simultaneously, and neither is a $C_0$ separation required together
with an $N_i$ separation.
\end{lemma}

\begin{proof}
The four increases of framing in the positive path replace the
$(-2)$ trace by a $2$ trace blown up at four distinct points of its
core.  This is the corner-blowdown description in
Proposition~\ref{prop:stack:lattice}.  Its cocore disk can be chosen
away from those four points.  On adjoining the reverse $2$ trace, the
two cocores join to a sphere.  The two normal disk trivializations
agree along their common boundary with opposite boundary orientations,
so the normal Euler number of this double is zero.  The double meets
each neighboring end core once.  With the orientations
\eqref{eq:geometry:oriented-spheres}, these intersections are positive.
The cocore double lies between the two middle cuts, whereas $X_i$
passes through its own middle cut $J_i$.  Consequently the double is
disjoint from both $J_i$.

At each intersection choose oriented complex disk coordinates in
which the two spheres are the coordinate axes.  Plumbing their normal
disk bundles by exchanging base and fiber disks gives a neighborhood
of the embedded configuration; its Euler numbers are $-4,0,-4$.
Taking sufficiently small disk fibers gives $C_2$ inside the macro.
The subconfiguration $U\cup X_{3-i}$ misses $J_i$, so its smaller
regular neighborhood $C_{1,i}$ also misses $J_i$.  Shrink $N_{3-i}$
inside it and choose a still smaller $C_0$ common to both subchain
neighborhoods.  This proves the asserted geometric placements.

The restrictions on simultaneous interfaces follow from the parameter
regions used below: a low-side separation has the other coefficient
bounded away from zero in normalized units, and thus cannot coexist
with the other low-side separation.  An opposite high-side $N_{3-i}$
separation is inside $C_{1,i}$.  The region using only $C_0$ has both
coefficients small, so no high-side separation occurs there.  We
therefore need only nested subchain neighborhoods and disjoint disk
bundle neighborhoods; small isotopies make their collars disjoint.
\end{proof}

For $0<a,b_0<1/4$, consider the polygon
\begin{equation}
 \mathcal P(a,b_0)=
 \{(x,y):4|y|\le x\le1,\ -a\le y\le b_0\}.
 \label{eq:geometry:polygon}
\end{equation}
The angle lattice is $2\pi\Z^2$.  In cyclic order the inward normals
are
\begin{equation}
 n_p=(1,4),\quad n_1=(0,1),\quad n_o=(-1,0),\quad
 n_2=(0,-1),\quad n_q=(1,-4).
 \label{eq:geometry:normals}
\end{equation}
The vertex $(0,0)$ has determinant eight.  We always remove its cone
point and attach the resulting outer lens boundary to the macro
exterior.  Figure~\ref{fig:geometry:polygon} shows the compact divisors
that remain.

\begin{figure}[htbp]
\centering
\begin{tikzpicture}[scale=3.2,>=Latex]
 \fill[blue!4] (0,0)--(.48,-.12)--(1,-.12)--(1,.17)--(.68,.17)--cycle;
 \draw[thick] (.08,-.02)--(.48,-.12)--(1,-.12)--(1,.17)--(.68,.17)--(.08,.02);
 \draw[dashed] (.08,-.02) arc[start angle=-90,end angle=90,x radius=.018,y radius=.02];
 \node[left] at (.04,0) {$\partial C_2$};
 \node[below] at (.76,-.12) {$X_1$};
 \node[right] at (1,.025) {$U$};
 \node[above] at (.84,.17) {$X_2$};
 \node[below left] at (.29,-.075) {$x+4y=0$};
 \node[above left] at (.39,.10) {$x-4y=0$};
 \draw[->,gray] (1.24,-.12)--(1.24,.17);
 \node[right] at (1.24,.025) {$a+b_0$};
 \begin{scope}[shift={(2.03,.025)},xscale=.36,yscale=.36]
  \draw[thick] (-.75,0)--(.75,0);
  \foreach \x/\lab/\wt in {-.75/X_1/-4,0/U/0,.75/X_2/-4}{
   \fill[white] (\x,0) circle (.075);
   \draw[thick] (\x,0) circle (.075);
   \node[above] at (\x,.09) {$\lab$};
   \node[below] at (\x,-.09) {$\wt$};
  }
 \end{scope}
\end{tikzpicture}
\caption{The punctured polygon and its compact divisor chain.  The
three affine edge lengths are $1-4a$, $a+b_0$, and $1-4b_0$.  Shrinking
a horizontal edge separates its $(-4)$ disk bundle.  Sending $a$ or
$b_0$ to zero instead moves the outer cone point to a subchain
interface; the true middle-sphere cut is introduced outside that
interface.}
\label{fig:geometry:polygon}
\end{figure}

\begin{lemma}[The toric model and its periods]
\label{lem:geometry:toric}
The toric space associated to the punctured polygon is $C_2$.  With
$w=a+b_0$ and normalized coordinates
\begin{equation}
 \xi=x/w,\quad y'=y/w,\quad L=1/w,\quad
 r_1=a/w,\quad r_2=b_0/w,
 \qquad r_1+r_2=1,
 \label{eq:geometry:normalization}
\end{equation}
it has a toric K\"ahler metric
\begin{equation}
 \begin{split}
 g&=\tfrac12 M_{ij}\,d\mu_i\,d\mu_j
       +2(M^{-1})_{ij}\,d\theta_i\,d\theta_j,
       \qquad (\mu_1,\mu_2)=(\xi,y'),\\
 M&=C+\sum_{\nu}\frac{n_\nu n_\nu^t}{l_\nu},
       \qquad C=\begin{pmatrix}1&0\\0&0\end{pmatrix},\\
 (l_p,l_1,l_o,l_2,l_q)
   &=(\xi+4y',r_1+y',L-\xi,r_2-y',\xi-4y').
 \end{split}
 \label{eq:geometry:hessian}
\end{equation}
Its positive K\"ahler ray is \eqref{eq:geometry:period}, and
\begin{equation}
 HX_1=1-4a,\quad HU=a+b_0,\quad HX_2=1-4b_0,
 \qquad H^2=2(a+b_0)-4(a^2+b_0^2)>0.
 \label{eq:geometry:period-pairings}
\end{equation}
At $a=0$ or $b_0=0$, with $w>0$, the retained piece is the indicated
punctured subplumbing $C_{1,i}$.  At a coefficient $1/4$, the retained
piece is punctured at the collapsed $X_i$ divisor.  The affine term
for a collapsed facet is retained in the limiting potential.  These
limits may occur together.
\end{lemma}

\begin{proof}
Every successive normal pair except $(n_q,n_p)$ has determinant one.
The corresponding torus corner chart is a smooth bidisk: the two
circle subgroups normal to its sides collapse on its coordinate axes.
Gluing these charts along an edge identifies the normal disk fibers
with the framing determined by the third normal.  The relations
\[
 n_p+n_o=4n_1,\qquad n_1+n_2=0,\qquad n_o+n_q=4n_2
\]
therefore give Euler numbers $-4,0,-4$, respectively.  At the remaining
corner $|\det(n_q,n_p)|=8$, so its removed neighborhood is a cone on
a lens space of order eight.  These charts construct exactly the
plumbing in Lemma~\ref{lem:geometry:plumbing}, without changing its
underlying smooth four-manifold.

The matrix $M/2$ is the Hessian of
\[
 \tfrac14\xi^2+\tfrac12\sum_\nu l_\nu\log l_\nu.
\]
It is positive definite in the polygon interior since its normals
span $\R^2$.  The coefficient $\tfrac12$ of each primitive boundary
term gives the smooth polar-coordinate extension along a simple
facet.  At a determinant-one corner this holds in both disk
coordinates; at the removed corner it holds on the finite quotient
chart.  This is the toric boundary convention of
\cite{Guillemin1994}, also verified directly by setting the normal
moment variable equal to half the square of a radial coordinate.

For the unnormalized polygon the affine lengths of the compact
facets are $1-4a$, $w$, and $1-4b_0$; symplectic areas are $2\pi$
times these lengths.  The intersection matrix in the ordered basis
$(X_1,U,X_2)$ is
\begin{equation}
 \begin{pmatrix}-4&1&0\\1&0&1\\0&1&-4\end{pmatrix}.
 \label{eq:geometry:plumbing-matrix}
\end{equation}
It is nonsingular with determinant eight and signature $(1,2)$.
Solving for the class with these three periods gives
$U+aX_1+b_0X_2$.  Normalizing the metric divides its K\"ahler form by
$w$ and hence does not change this ray.  Direct multiplication gives
\eqref{eq:geometry:period-pairings}; positivity follows also from
$a^2+b_0^2\le(a+b_0)/4$.

If $a=0$, the $X_1$ coefficient is absent, so the class is carried by
the subchain $U,X_2$.  The new vertex is removed, while the affine
term belonging to the collapsed outer edge remains as a redundant
term in the Hessian.  The other low limit is identical with the
indices exchanged.  If $a=1/4$, then $HX_1=0$.  The rational
orthogonal decomposition at $\partial N_1$ consequently puts the
whole class on the complement of $N_1$.  The same argument applies
to $X_2$.  The subchain matrix has determinant $-1$; each isolated
disk bundle has determinant $-4$.  Thus all these separating
boundaries are rational homology spheres, so the decompositions over
$\R$ are nondegenerate.  Removing only negative directions, or
retaining the subchain of signature $(1,1)$, leaves positive rank one.
\end{proof}

\subsection{Scalar curvature of the Hessian metrics}

We give the curvature calculation because its quantitative strictness
is needed when the tips are changed conformally.  It applies to both
the main polygon and the satellite metrics.

\begin{lemma}[Curvature and dependence bounds]
\label{lem:geometry:scalar}
Let $l_\nu$ be positive affine functions with gradients $n_\nu$
spanning $\R^2$, and let $C\ge0$ be constant.  For the metric
\eqref{eq:geometry:hessian}, put
\[
 B=M^{-1},\quad
 A_\nu=B^{1/2}n_\nu/\sqrt{l_\nu},\quad
 P_{\nu\mu}=A_\nu^tA_\mu,\quad
 z_\nu=l_\nu^{-1/2},\quad q_\nu=z_\nu P_{\nu\nu}.
\]
Then $0\le P\le I$, $\Scal(g)=2\mathcal S$, and
\begin{align}
 \mathcal S
 &=2\sum_\nu z_\nu^2P_{\nu\nu}^2
     -q^tPq-z^tP^{\circ3}z \notag\\
 &=q^t(I-P)q+
 z^t\bigl((I-P)\circ(P\circ P)\bigr)z\ge0.
 \label{eq:geometry:scalar-identity}
\end{align}
Here $\circ$ is entrywise product.  For every nonzero relation
$\sum_\nu\gamma_\nu n_\nu=0$, its two nonnegative terms satisfy
\begin{align}
 q^t(I-P)q
 &\ge \frac{(\sum_\nu\gamma_\nu P_{\nu\nu})^2}
                    {\sum_\nu\gamma_\nu^2l_\nu},
 \label{eq:geometry:first-bound}\\
 z^t\bigl((I-P)\circ(P\circ P)\bigr)z
 &\ge\frac{\left\|\sum_\nu\gamma_\nu A_\nu A_\nu^t\right\|_{\mathrm{HS}}^2}
                   {\sum_\nu\gamma_\nu^2l_\nu}.
 \label{eq:geometry:second-bound}
\end{align}
The main metrics have strictly positive scalar curvature on compact
sets away from punctures, including their smooth toric axes.  This
strictness remains uniform on fixed compact truncations of their
$L\to\infty$ tips and of the three-direction satellites below.
\end{lemma}

\begin{proof}
Let $A$ be the matrix with columns $A_\nu$.  Then
\[
 AA^t=B^{1/2}(M-C)B^{1/2}
       =I-B^{1/2}CB^{1/2}\le I.
\]
Hence $P=A^tA$ is a positive contraction.  For completeness set
$G_{\nu\mu}=n_\nu^tBn_\mu$.  Differentiating
$B M=I$ gives
\[
 \partial_kB
   =\sum_\nu \frac{n_{\nu k}}{l_\nu^2}
                     Bn_\nu n_\nu^tB.
\]
Differentiating once more, contracting the two derivative indices,
and collecting the derivative of $l_\nu^{-2}$ separately from the
two derivatives of $B$ yields
\[
 -\partial_i\partial_jB_{ij}
 =2\sum_\nu\frac{G_{\nu\nu}^2}{l_\nu^3}
 -\sum_{\nu,\mu}
   \frac{G_{\nu\nu}G_{\mu\mu}G_{\nu\mu}+G_{\nu\mu}^3}
        {l_\nu^2l_\mu^2}.
\]
The real Riemannian scalar curvature of a toric metric is minus the
contracted second derivative of its inverse potential Hessian
\cite[Theorem~4.1(i), arXiv version~1]{Abreu1998}; it is twice the K\"ahler scalar
curvature denoted by $R$ in that reference.
That inverse Hessian is $2B$ here, proving
\eqref{eq:geometry:scalar-identity}, including its factor two.  The
second equality uses
$\operatorname{diag}(P_{\nu\nu}^2)-P^{\circ3}
=(I-P)\circ(P\circ P)$.  Both summands are nonnegative by the
Schur product theorem.

Let $k_\nu=\gamma_\nu\sqrt{l_\nu}$.  The relation on the normals
gives $Ak=0$, hence $Pk=0$ and
\[
 I-P\ge \frac{kk^t}{\|k\|^2}.
\]
Pairing with $q$ proves \eqref{eq:geometry:first-bound}.
Schur-multiplying the same inequality by $P\circ P$ and pairing
with $z$ proves \eqref{eq:geometry:second-bound}, because
\[
 \sum_{\nu,\mu}\gamma_\nu\gamma_\mu
                 (A_\nu^tA_\mu)^2
 =\left\|\sum_\nu\gamma_\nu A_\nu A_\nu^t\right\|_{\mathrm{HS}}^2.
\]

For the full polygon use the positive dependence
$\gamma=(1,1,2,1,1)$ in the order $(p,1,o,2,q)$.
Its first numerator is strictly positive in the interior.  At a
simple facet its active column $A_\nu$ tends to a unit vector, so
$P_{\nu\nu}\to1$ and this numerator remains positive.  At a smooth
vertex the two active columns tend to orthonormal vectors, with the
same conclusion.  These assertions follow either by inverting $M$
in active-normal coordinates or by taking the Schur complement of
its singular normal term.  They also prove a positive lower bound
on any compact smooth truncation.

For the tip with $\xi=O(1)$ and $L\to\infty$, discard the vanishing
$n_o$ term and use $n_1+n_2=0$.  Its denominator is
$l_1+l_2=1$, and its first numerator is
$(P_{11}+P_{22})^2$.  In the interior this is positive.  On an
active horizontal edge one diagonal term tends to one.  On an
active sloping edge, the projections of $n_1,n_2$ to the remaining
tangent direction are nonzero; inversion on that direction again
gives a positive limit.  At a smooth vertex one horizontal edge is
active.  Thus strictness extends to all smooth boundary points of
each fixed truncated tip.

After satellite rescaling the constant $C$ disappears and exactly
three distinct normal directions remain.  Their unique relation
has three nonzero coefficients.  The three rank-one symmetric
matrices $n_\nu n_\nu^t$ are linearly independent: in coordinates
where two directions are the axes, the third has a nonzero
off-diagonal entry.  Conjugation by $B^{1/2}$ and multiplication by
positive $l_\nu^{-1}$ preserve independence.  Thus the numerator
of \eqref{eq:geometry:second-bound} is nonzero in the interior.
At a simple edge its active unit column is orthogonal in the limit
to the other columns, so the corresponding nonzero coefficient
prevents cancellation.  At a smooth vertex the two active columns
are orthonormal and again prevent cancellation.  The same bound
is positive on each compact truncation, including the smooth axes.
\end{proof}

\begin{lemma}[Cone links]
\label{lem:geometry:cones}
At every puncture, the homogeneous Hessian formed from precisely the
vanishing linear terms is a metric cone.  Its smooth link has scalar
curvature at least six.  The link at a collapsed $X_i$ admits a path
through positive-scalar-curvature metrics to the round lens metric.
\end{lemma}

\begin{proof}
Write $E=\mu_i\partial_{\mu_i}$ for the moment Euler field at the
vertex.  For homogeneous $l_\nu$ and $C=0$,
\[
 g(E,\cdot)=\tfrac12 d\sum_\nu l_\nu,\qquad
 R^2=2\sum_\nu l_\nu.
\]
The dilation of moment variables by $s^2$ multiplies $g$ by $s^2$,
and $2E=R\partial_R$ has squared length $R^2$.  Thus
$g=dR^2+R^2h$.  Its scalar curvature is
$R^{-2}(\Scal(h)-6)$, so Lemma~\ref{lem:geometry:scalar} gives
$\Scal(h)\ge6$.  Each bounding side has primitive normal.  The
torus action therefore has no finite isotropy along its collapsed
circle, and the link is smooth, even when the cone vertex itself
is a finite quotient singularity.

At $N_1$ the adjacent linear forms are $p$ and $v$ below, and the
redundant form is $(p+v)/4$.  Multiply its Hessian contribution by
$s\in[0,1]$.  For $s>0$ this is the contribution of the positive
linear form $s(p+v)/4$ itself, since
$\nabla(sl)\nabla(sl)^t/(sl)=s\nabla l\nabla l^t/l$.
The nonnegative scalar calculation applies for every $s$, including
the two-form limit $s=0$.  Identifying the radius-one links by
radial dilation gives a smooth path with scalar curvature at least
six.  At $s=0$ the two independent primitive-side terms give the
flat quotient of $\C^2$, whose link is the round lens space.  The
other side is obtained by reflection.
\end{proof}

\subsection{Shrinking edges, satellites, and cylindrical completions}

In the normalized polygon the parameters measuring an edge shrink
are
\begin{equation}
 r_i,\qquad h_i=L/4-r_i.
 \label{eq:geometry:shrink-parameters}
\end{equation}
Only one $r_i$ can be small, because $r_1+r_2=1$.  A small $h_i$
occurs near the corresponding right-hand corner and separates
$N_i$.  A small $r_i$ occurs at a different corner and separates
$C_{1,i}$.  These zones can be chosen disjoint.  The sole nesting
is the original determinant-eight puncture inside an $r_i$
satellite.  We retain this distinction in choosing the conformal
factors.

\begin{lemma}[Uniform collar and satellite models]
\label{lem:geometry:collars}
For every fixed derivative order $k$, the metric near an edge with
shrink parameter $\delta=r_i$ or $h_i$ has the following description.
There is a homogeneous radius $R$ and a cone metric $g_{\mathrm c}$
such that, in scale-invariant $C^k$ norms on every dyadic annulus,
\begin{equation}
 \|g-g_{\mathrm c}\|_{C^k(g_{\mathrm c})}
 \le C_k\bigl(\delta/R^2+R^2\bigr),
 \qquad C\sqrt\delta\le R\le R_*.
 \label{eq:geometry:cone-error}
\end{equation}
The estimates include smooth toric axes and are uniform on the
$L\to\infty$ tips.  In affine moment coordinates $\mu_{\mathrm{loc}}$
centered at the relevant homogeneous cone vertex, the scale
$\mu_{\mathrm{loc}}=\delta\widehat\mu_{\mathrm{loc}}$ gives
$g=\delta(g_{\mathrm{sat}}+o(1))$ on compact sets away from any
remaining cone point.  The metric $g_{\mathrm{sat}}$ is a
three-direction Hessian metric with strictly positive scalar
curvature on those compact sets.

There are conformal changes making all the necessary annuli
asymptotically cylindrical, with fixed positive limiting link
sizes.  They have nonnegative scalar curvature and preserve
strict positivity on smooth compact regions.  Arbitrarily deep
in those annuli one may interpolate to the exact limiting product
over a bounded logarithmic width, with $C^k$ errors tending to
zero.  In particular the separated $N_i$ cap can be made positive
everywhere, including its exact product end.
\end{lemma}

\begin{proof}
The two local calculations are explicit.  At a small $h_1$, set
\[
 p=\xi+4y',\qquad v=L-\xi.
\]
The three relevant forms are $p,v,(p+v)/4-h_1$.  Their homogeneous
limits are $p,v,(p+v)/4$.  At a small $r_1$, set
\[
 q=\xi-4y',\qquad l=r_1+y'.
\]
The relevant forms are $q,l,q+8l-8r_1$, with homogeneous limits
$q,l,q+8l$.  Reflection gives side two.  In each case use the sum
of the three homogeneous forms to define $R$ as in
Lemma~\ref{lem:geometry:cones}.

Only the redundant homogeneous form is shifted.  It is uniformly
comparable to $R^2$: in the high model $(p+v)/4=R^2/10$, and in
the low model $R^2/4\le q+8l\le4R^2/9$.  Replacing that
denominator by its shift of size $O(\delta)$ therefore gives a
relative error $O(\delta/R^2)$.  The bounding side forms are
unchanged, including where one vanishes on an axis.  The constant
matrix and the remaining regular
facet terms give a relative error $O(R^2)$.  Differentiating after
rescaling the annulus to unit size gives the same bound at every
fixed derivative order.  Near a simple axis, use the square root
of its vanishing affine form as the radial disk coordinate.  The
singular normal term is then precisely the smooth polar disk
metric, and its coefficients and those of the inverse Hessian
obey the same estimates.  This proves
\eqref{eq:geometry:cone-error} through the axis.  At the left tip
as $L\to\infty$, the omitted $l_o^{-1}$ term is bounded with all
scaled derivatives and tends to zero; all other omitted
denominators stay bounded away from zero on the chosen zone.
Thus the constants are uniform there.

Use $\mu_{\mathrm{loc}}=(p,v)$ at a high-side shrink and
$\mu_{\mathrm{loc}}=(q,l)$ at a low-side shrink.  In the variables
$\widehat\mu_{\mathrm{loc}}=\mu_{\mathrm{loc}}/\delta$, the three
singular Hessian terms have factor $\delta^{-1}$ while
$d\mu_{\mathrm{loc}}$ has factor $\delta$.
The moment and angle terms of the metric both have
factor $\delta$.  Removing that factor leaves exactly the three
affine-direction metric.  The constant and regular terms tend
to zero.  Lemma~\ref{lem:geometry:scalar} proves its asserted
strict positivity.

Choose a nondecreasing smooth function $\chi_\infty$ which is zero
below a large fixed number and one above a larger fixed number,
and put
\begin{equation}
 F(t)=\int_t^\infty\chi_\infty(s)s^{-2}\,ds,
 \qquad
 f=1+c\,\chi_0(R)\delta^{-1/2}F(R/\sqrt\delta).
 \label{eq:geometry:conformal-factor}
\end{equation}
Here $\chi_0$ is one at small fixed radius and cuts off at a
slightly larger fixed radius.  Choose those radii inside the
zone above.  The new metric is $f^2g$.  In dimension four its
scalar curvature is $f^{-3}(-6\Delta_g+\Scal(g))f$, where
$\Delta_g=\operatorname{div}\grad$.

For $q(R)=\delta^{-1/2}F(R/\sqrt\delta)$ a direct calculation gives
\begin{equation}
 \Delta_g q
 =-\frac{\chi_\infty'(R/\sqrt\delta)|dR|^2}
          {\sqrt\delta\,R^2}
   +\frac{\chi_\infty(R/\sqrt\delta)}{R^3}
       \bigl(2|dR|^2-R\Delta_gR\bigr).
 \label{eq:geometry:radial-laplacian}
\end{equation}
For the exact cone, $R\Delta R=3|dR|^2$.  By
\eqref{eq:geometry:cone-error}, choosing the initial ramp point
large and $R_*$ small gives
$R\Delta_gR>(2+\sigma)|dR|^2$ for one fixed $\sigma>0$
where $q$ starts decreasing.  Thus $-\Delta_gq$ is positive
there.  Where $q$ is constant, strict positivity of the base
scalar curvature suffices.  On the outer cutoff annulus, the
main or rescaled-tip metrics have a common positive scalar
lower bound; choosing $c>0$ sufficiently small absorbs all
derivatives of $\chi_0$.  We obtain
$(-6\Delta_g+\Scal(g))f>0$ on the smooth part.  The coefficient
$c$ may be tapered smoothly to zero when $\delta$ leaves its
small-parameter zone.  Choose the taper only where the same
strict compact bounds hold.  The limiting factor is $1+c/R$
on the cylindrical annulus, with fixed $c>0$ near $\delta=0$.

At the original puncture the homogeneous radius is
$R_8=\sqrt{4\xi}$.  Add a cutoff multiple $c_8/R_8$.  If no
$r_i$ is small, this is the preceding construction without an
inner smoothing scale.  If $r_i\to0$, place its cutoff at a
fixed sufficiently small radius in the rescaled satellite,
that is, at $R_8$ comparable to a small fixed multiple of
$\sqrt{r_i}$.  It then lies where the first factor in
\eqref{eq:geometry:conformal-factor} is constant of order
$r_i^{-1/2}$.  Write $R_8=\sqrt{r_i}\rho$ and cancel this scale
against $g=r_i g_{\mathrm{sat}}+o(r_i)$.  The added factor is
a fixed multiple of $c_8/\rho$ on a fixed-scale satellite.
Its cutoff error is therefore absorbed by the satellite's
strict compact positivity after choosing $c_8$ small.  This
also shows that the nested link does not shrink to zero.
At most one such nesting occurs.  The other zones are
disjoint, so their choices do not interfere.

On a cylinder use $s=-\log R$.  The cone metric multiplied by
$(c/R)^2$ is $c^2(ds^2+h)$.  The difference between $1+c/R$
and $c/R$ has relative size $O(R/c)$, in addition to the error
\eqref{eq:geometry:cone-error}.  Interpolating to
$c^2(ds^2+h)$ over a fixed $s$-width deep in the annulus has
errors tending to zero with all fixed spatial derivatives.
Because the product scalar curvature is bounded positively
away from zero, a sufficiently deep interpolation preserves
positivity.  The same argument applies at the nested puncture
after rescaling.  On a separated high-side satellite, the
compact core and the entire outward conformal annulus are
positive; the exact product continuation is positive as well.
This is the required complete positive metric on $N_i$.
\end{proof}

\subsection{The interval cube and its true ends}

Use coordinates $t_i\in[-3,3]$ on the two intervals of a macro.
The true $J_i$ face is $t_i=-3$, and the true $N_i$ face is
$t_i=3$.  Choose a smooth nondecreasing function $\alpha_i(t_i)$
which is zero for $t_i\le0$, equals $1/4$ for $t_i\ge1$, and
increases between these values on $(0,1)$.  It may be taken
stationary to all orders at the endpoints.  A sphere test keeps
its whole-loop rule for $t_i\ge0$.  Its conversion to the two
segment rules of Section~\ref{sec:stack} takes place only at
$t_i<0$, after that side is no longer used as a whole test.

Set $w_0=\alpha_1+\alpha_2$.  For $w_0\ge3\epsilon$ the
target coefficients are $(\alpha_1,\alpha_2)$.  For a finite
smooth metric replace them by
\begin{equation}
 a_i=\alpha_i+d w_0(1-8\alpha_i),\qquad
 (a_1,a_2)=(a,b_0),\qquad
 w=w_0\bigl(1+d(2-8w_0)\bigr).
 \label{eq:geometry:regularization}
\end{equation}
For a common sufficiently small $d>0$, both coefficients lie
strictly between zero and $1/4$.  At low saturation
$\alpha_i=0$, and at high saturation $\alpha_i=1/4$, respectively,
the shrink parameter is
\begin{equation}
 r_i=\frac{d}{1+d(2-8w_0)},\qquad
 h_i=\frac{d}{1+d(2-8w_0)}.
 \label{eq:geometry:saturation-scale}
\end{equation}
Thus there is one uniform small scale for productizing at either
saturation.

\begin{proposition}[Assembly in the positive-width region]
\label{prop:geometry:assembly}
For $w_0\ge3\epsilon$, the toric models and their collars form
a smooth family on the fixed macro.  The low and high saturation
regions extend to the prescribed true faces.  At fixed $\epsilon$,
as $d\to0$, all retained positive-width main pieces converge
on compact subsets to the complete conformal metrics of
Theorem~\ref{thm:geometry:family}\textup{(ii)}, uniformly under
further genuine length insertions.
\end{proposition}

\begin{proof}
Always remove the original outer cone point, complete it by
Lemma~\ref{lem:geometry:collars}, and attach its product collar
to the exterior of $C_2$.  The collar is an artificial finite
collar.  Its productization may be moved arbitrarily far down
the puncture as $d\to0$, so on every compact subset the
retained metric is exactly conformal to
\eqref{eq:geometry:hessian} in the limit.

Near a saturation, productize at $R\asymp d^{1/4}$ whenever
the corresponding $\delta$ is at most $2d$, and turn off this
productization by $\delta=3d$.  Equations
\eqref{eq:geometry:cone-error} and
\eqref{eq:geometry:saturation-scale} show that its Hessian error
is $O_k(\sqrt d)$, while its conformal product error is
$O_k(d^{1/4})$.  A fixed-width cutoff in $\log R$ therefore
has vanishing errors.  Use the adjacent affine forms to fix
the angular coordinates of this product.  They determine its
link independently of all tangential parameters.

For $t_i\ge1$, this gives an exact collar separating $N_i$.
Attach any additional length inside the exact product, with
length zero initially and tending to infinity only as
$t_i\to3$.  Keep the link and the coefficient $c$ fixed on
that collar.  The satellite on the $N_i$ side is the positive
cap of Lemma~\ref{lem:geometry:collars}.  Its metric, like
every retained sector, is unaffected by length inserted
in another disjoint product collar.  The metric on the
complementary sector is unchanged by the bit supported
on $X_i$.

For $t_i\le0$, the collar instead separates $C_{1,i}$.  It
must be retained while the exterior is changed to introduce
$J_i$.  Since $w_0\ge3\epsilon$, the other coefficient is
positive; there is no competing second low separation.
First extend the outer-puncture productization back to a
fixed sufficiently far-out portion of the rescaled satellite,
where it is already close to a positive product.  Shorten
that obsolete outer collar there.  Similarly extend the
product structure a bounded distance toward the compact
part on the $C_{1,i}$ side, while retaining the auxiliary
$\partial C_{1,i}$ collar, whose length tends to infinity as
$d\to0$.  For fixed auxiliary choices this collar remains finite
as $t_i\to-3$.  All of these extensions
take place in positive-product regions, so they have bounded
geometry and preserve positive scalar curvature along their
long parts.

The exterior of this retained collar now consists of bounded
pieces.  On those pieces interpolate to a metric with a
round product collar at $J_i$, and fixed data on its far
side.  This is possible because $C_{1,i}$ misses $J_i$ by
Lemma~\ref{lem:geometry:plumbing}.  Keep the inner collar
fixed during the interpolation.  Only after reaching the
round product data insert the genuine $J_i$ length, tending
to infinity as $t_i\to-3$.  The far side was made independent
of the other interval parameter before this insertion.
Consequently the true split has the required local parameter
dependence.  The retained main metric and its long
$\partial C_{1,i}$ collar have not been changed.

Use successive closed subintervals of the low or high
saturation range for these operations, with each operation
stationary near the joints.  If another parameter has no
active operation it simply waits.  Disjoint genuine collars
then permit independent insertion lengths.  In a low-high
corner the $N_{3-i}$ collar is inside $C_{1,i}$ and its
length insertion commutes with the exterior low-side
operation.  This gives the asserted compatibility.

It remains to justify that these varying polygons and
interfaces are placed on one fixed manifold.  Start from a
fixed punctured plumbing polygon, with marked cross curves
for the interfaces actually in use.  A cross curve near a
shrinking $X_i$ edge bounds that divisor's disk bundle.  A
cross curve at small $r_i$ joins the $X_i$ edge to the far
sloping side; its inner side contains exactly the subchain
$U,X_{3-i}$.  A cross curve joining the two horizontal sides
in the long strip bounds the chosen $U$ neighborhood.  These
are exactly the neighborhood types in
Lemma~\ref{lem:geometry:plumbing}.

On each side of the polygon, place the endpoints of these
curves by increasing coordinate changes.  Their order does
not change, since only compatible disjoint or nested curves
are marked together.  Choose the curves as segments or
graphs in rectangular side charts and extend their motion
by vector fields tangent to the polygon faces, supported
away from unmarked corners.  On the remaining disk interiors
extend the motion with a cutoff.  This constructs smooth
face-preserving diffeomorphisms carrying the marked collars
to the fixed interfaces.  Near a face they are smooth
positive rescalings of its normal variable, so their lifts
in the square-root disk coordinates are smooth.  Carry the
angle torus along with these maps; its lattice and all
transition functions are fixed.  The two low placements use
disjoint parameter zones and fix the common smaller
$C_0$ whenever it is marked.  Thus they impose no crossing
condition on one another.

For any exterior interpolation first shorten unretained
long products, trivialize the resulting bounded pieces by
these maps, and interpolate the positive-definite tensors
there with their product collars fixed.  Reinsert retained
products afterward.  This is a convex interpolation on a
compact family of bounded pieces and hence gives uniform
spatial geometry; it does not require a uniform parameter
derivative in a stretching coordinate.  The construction
is stationary on smaller genuine face collars except for
the prescribed lengths and retained tangential parameters.
It proves both smooth assembly and the stated convergence.
\end{proof}

\subsection{Small width and the long spherical tube}

\begin{lemma}[The tube and the corner replacement]
\label{lem:geometry:tube}
The family of Proposition~\ref{prop:geometry:assembly} extends
across $w_0\le3\epsilon$ to the low-low corner, with exactly
the true cuts $J_1,J_2$.  Every macro in this region contains
a tube of fiber class $U$ and length comparable to
$\epsilon^{-1}$.  The extension satisfies the uniform
description in Theorem~\ref{thm:geometry:family}\textup{(iii)--(iv)}.
\end{lemma}

\begin{proof}
First suppose $w_0\in[\epsilon,3\epsilon]$, so
$L\asymp\epsilon^{-1}$ and no high coefficient is saturated.
Put $z=y'+r_1$.  On the strip between the two horizontal
faces, $0\le z\le1$.  The two horizontal Hessian terms add
to $1/[z(1-z)]$ in the $y'$ direction.  If both $\xi$ and
$L-\xi$ tend to infinity, the other facet terms tend to zero.
Thus the limiting metric is
\begin{equation}
 \frac{dz^2}{2z(1-z)}+2z(1-z)d\theta_y^2
       +\tfrac12d\xi^2+2d\theta_x^2.
 \label{eq:geometry:explicit-tube}
\end{equation}
Under $z=(1+\cos\varphi)/2$ its first two terms are
$\tfrac12(d\varphi^2+\sin^2\varphi\,d\theta_y^2)$.
This proves the radius and normalization in
\eqref{eq:geometry:long-model}.  The class of the spherical
fiber is $U$, since it is the toric sphere transverse to
the longitudinal strip between the two horizontal facets.

For every fixed derivative order, on
$T\le\xi\le L-T$ the error is bounded by a number
$\eta_k(T)\to0$ as $T\to\infty$, uniformly in $r_i\in[0,1]$.
This follows by differentiating the remaining reciprocal
affine terms; the denominators are at least a fixed
multiple of the distance to a strip end when $T$ is large.
Using $\sqrt z$ and $\sqrt{1-z}$ near the two sphere poles
gives the same conclusion through the axes.  All conformal
tip changes were supported in bounded end regions and
do not meet this bulk strip.

Choose a large fixed $T$.  First interpolate on the strip
to make the metric exactly \eqref{eq:geometry:long-model}
between $2T$ and $L-2T$, using cutoffs in fixed-width
transition strips.  The interpolated errors obey the same
bulk bounds, in particular with two derivatives.  Within the parameter
region $\epsilon\le w_0\le3\epsilon$, deform the toric model from regularizer $d$
to one fixed small positive value.  The obsolete
artificial collar lengths then become bounded.  Throughout
these stages retain the low-saturation exterior procedure
of Proposition~\ref{prop:geometry:assembly} if either
parameter is low.  It changes only an exterior of the
retained subchain and does not reach the tube.

Now bring the exact tube to a fixed length comparable to
$\epsilon^{-1}$.  The remaining tips, in normalized
coordinates, have bounded size, and their effective
$r_i$ are bounded away from zero because the regularizer
has been fixed positively.  Interpolate these bounded
pieces to template metrics with product collars at both
$J_i$, retaining the exact tube.  Its placement inside
$C_0$ is disjoint from both $J_i$.  If a $J_i$ insertion
has already started, perform the interpolation on the
underlying shortened metric with its product collar
fixed, then reinsert precisely the same length.  Arrange
the original low-side stages to supply this product
before the corresponding insertion starts.  The far
side stays fixed, so no dependence on the other interval
parameter is introduced at a true cut.

These successive operations can be assigned disjoint
subranges of $w_0/\epsilon$, with stationary overlaps.
Below $w_0=\epsilon$ retain only the fixed long tube,
the bounded templates, and the ordinary low-end
prescriptions.  This includes $w_0=0$, where neither a
normalized polygon nor a positive-width period is used.
Both genuine $J_i$ lengths are still inserted in their
own collars.  Thus the corner has no extra face or
extra multiplicity.

For the uniform assertion, take any sequence of parameters.
The identities $r_1+r_2=1$ and $L\ge2$ imply that the two
small-$r_i$ parameter regimes are mutually exclusive.
Also $h_i\ge1/2-r_i$ excludes simultaneous small $r_i,h_i$
on the same side.  A low degeneration and the opposite
high degeneration, or the two high degenerations, may
occur simultaneously; their spatial neighborhoods are
disjoint.  The only nested spatial degeneration is the
already described outer puncture inside a low satellite.
Lemma~\ref{lem:geometry:collars} lists
the smooth main scale, cone annulus, and rescaled satellite
for each shrinking edge.  The only remaining possible
unbounded toric direction is $L\to\infty$, just analyzed
by \eqref{eq:geometry:explicit-tube}.  A deactivated
rational collar in a low homotopy was shortened inside
a positive product.  Every free exterior interpolation
was on bounded pieces.  Hence this list exhausts all
unbounded regions after passing to a subsequence.

For a toric macro retained above $3\epsilon$,
$L\le C/\epsilon$ by \eqref{eq:geometry:regularization};
in the replacement region its chosen length has the
same bound.  Fixing $T$ leaves a bounded number of
truncated tips, whose geometric costs are bounded by
a constant $C_T$, independent of deeper regularizations
and longer rational products.  A bounded local error
tending to zero on the tube bulk has integrated cost
at most
\begin{equation}
 C_T+\eta(T)\,C\epsilon^{-1},\qquad \eta(T)\longrightarrow0.
 \label{eq:geometry:integrated-error}
\end{equation}
Multiplying by $\epsilon$, letting $\epsilon\to0$ first,
and then $T\to\infty$ proves the stated
$o(\epsilon^{-1})$ assertion.  The rational product collars inside a macro
have positive scalar curvature and their own exact
models, so their arbitrary lengths contribute no error
of this type.  Their nonexact annular ends have the
same uniform bounded-truncation description.  This
establishes the full sequential assertion.
\end{proof}

Intervals outside positive macros may use any smooth interval
path with the same true cuts, taking the positive complete
disk-bundle metric of Lemma~\ref{lem:geometry:collars} at
the lens end and the round middle cut at the other end.
The positive-definite tensors with fixed collars form a
convex space on the shortened compact piece, so such paths
exist with stationary endpoints.  Different macro procedures
are products on disjoint regions.  The external rational
$Y_j$ collars and the old-cap exceptional-sphere isolation
collars are disjoint from all these modifications.  They
can therefore be inserted with arbitrary finite lengths.
This proves the metric and true-face parts of
Theorem~\ref{thm:geometry:family}.

\subsection{Cohomology support and uniform comparison forms}

We finish with the two topological interfaces needed in the
Abelian estimates.  They prevent a period comparison from
using a class on a piece which has disappeared, and prevent
long rational collars from adding an uncontrolled background
energy cost.

\begin{lemma}[Support at period vertices]
\label{lem:geometry:period-support}
Let $H$ be a positive-width limiting period in
\eqref{eq:geometry:period}.  Express it as a convex
combination of the four vertex classes
\[
 U,\qquad U+X_1/4,\qquad U+X_2/4,
 \qquad U+(X_1+X_2)/4.
\]
Every vertex with nonzero weight is carried rationally by
the retained smooth main piece.  It omits every side whose
coefficient is zero, and is orthogonal to every divisor
whose coefficient is $1/4$.  Its pairings may consequently
be computed with compactly supported Poincar\'e forms on
that piece, independently of any reference filling at a
true lens cut.  A side occurring in such a nonempty vertex
still has its whole-sphere test in the approximating
positive-width family.
\end{lemma}

\begin{proof}
Set $p=4a$, $q=4b_0$.  The four weights are
$(1-p)(1-q)$, $p(1-q)$, $(1-p)q$, and $pq$.
If $p=0$, all nonzero-weight vertices omit $X_1$ and lie
in the subchain $U,X_2$.  If $p=1$, all such vertices
contain $X_1/4$, and their intersection with $X_1$ is
$1+X_1^2/4=0$.  The same statements hold for $q$ and
$X_2$.  At a rational boundary, Mayer--Vietoris over
$\R$ splits second cohomology, and the intersection
form splits orthogonally.  Thus absence of a coefficient
gives support on the retained subchain, and orthogonality
to an isolated negative divisor gives support on its
complement.  The relative-to-absolute cohomology map is
surjective in degree two because the boundary has
$H^1=H^2=0$ over $\R$.  Each supported class therefore
has a compactly supported representative on the completed
main piece.

A positive limiting coefficient has $t_i>0$ unless it
is on the high saturated side, where $t_i\ge1$.
For all sufficiently late approximating parameters
that side is in the whole-loop range.  If its $N_i$
has been cut off it is also, by convention, a whole
test.  No side in the low switching range is used.
\end{proof}

\begin{lemma}[Comparison forms, uniformly in the auxiliary lengths]
\label{lem:geometry:comparison}
Fix a finite set of real degree-two classes on the fixed
stack, or on its genuine-cut components with specified
bounded reference extensions.  On every family in
Theorem~\ref{thm:geometry:family} each such class $\Lambda$
has a closed imaginary-valued representative $b_*$, with
$b_*/(2\pi i)$ representing $\Lambda$, such that the
following hold.  On each long tube of fiber $U$ it is
\begin{equation}
 b_*=i\lambda_U\,\operatorname{vol}_{S^2(1/\sqrt2)},
 \qquad \lambda_U=\Lambda U,
 \label{eq:geometry:comparison-flux}
\end{equation}
apart from bounded end transitions.  It vanishes on the
long middle of every rational collar.  For every fixed
truncation size $T$, its energy off the bulk tubes is
bounded by a constant $C_T$ on all compact and rescaled
pieces, independently of deeper regularizations and
longer rational collars.  On the tube bulk its norm and
volume differ from their product values by an error of
the form \eqref{eq:geometry:integrated-error}.  Consequently
\begin{equation}
 \int |b_*|^2=O(\epsilon^{-1}),\qquad
 \int_{\mathrm{off\ tubes}}|b_*|^2=o(\epsilon^{-1}),
 \label{eq:geometry:comparison-energy}
\end{equation}
where bounded end regions may be included in either
term.  The same uniform $o(\epsilon^{-1})$ assertion
holds for the squared negative scalar-curvature costs
on the region obtained by omitting the $Y_j$ middles
and retaining prefixes of length $O(\epsilon^{-1/2})$,
and for bounded local defects from the bulk model.
There is no scalar-curvature assertion on the omitted
$Y_j$ middles.  Constants may depend on
the fixed classes and finite problem.  These assertions
do not assert uniform bounds as the integral class
itself ranges over an unbounded set.
\end{lemma}

\begin{proof}
First shorten every long rational collar and every
long tube to bounded length, keeping its product
identification and the corresponding collapse map.
For a shrinking edge also cut in the cone annulus,
retain fixed compact truncations of its main and
rescaled satellite pieces, and shorten the intervening
product.  Do this inside the satellite for its nested
outer puncture if present.  The proof of
Lemma~\ref{lem:geometry:tube} shows that there are only
finitely many resulting decomposition types.  To check
compactness also when $L$ is unbounded, observe that then
both $h_i=L/4-r_i$ tend to infinity.  Off the bulk strip
there remain only the tip $\xi\le2T$, the tip
$L-\xi\le2T$, and possibly one low-side satellite.
At the latter, the rescaling and nested puncture
truncation of Lemma~\ref{lem:geometry:collars} give
fixed compact charts.  At the opposite tip use
$(L-\xi,z)$, where $z=y'+r_1$; the two sloping-facet
terms tend to zero, and the other terms converge
smoothly on the retained charts.  When $L$ is bounded,
ordinary parameter compactness together with the
finitely many edge rescalings gives the same conclusion.
The low-side exterior homotopies were performed on
bounded pieces.  Thus within each decomposition type
the truncated metrics and their collar identifications
vary in compact smooth families.
The moment sides and angle lattice are fixed; no
unbounded change of homology basis occurs in passing
to these truncations.  Uniform equivalence of these
metrics on each fixed truncation is enough for the
following de Rham construction.

Choose closed representatives of a basis for the
cohomology of each shortened manifold.  On a rational
product collar, a degree-two class restricts to zero
over $\R$.  Its representative there is therefore
$d\eta$, and subtracting $d(\chi\eta)$ sets it to
zero on a smaller middle collar.  On a tube slice,
$H^2(S^2\times S^1;\R)$ is generated by the spherical
area form.  Since that sphere has area $2\pi$,
\eqref{eq:geometry:comparison-flux} has exactly the
prescribed integral.  The difference between it and
the original representative is exact on a middle
slab.  The same cutoff construction replaces the
representative by the prescribed vertical form
there.  Use disjoint smaller collars for these
operations.  Local primitives can be chosen by a
fixed bounded right inverse on the shortened
collars, so their norms and the resulting energy
are bounded uniformly on each compact family.

Restore every long product using the same tangential
closed form on its whole middle: zero on a rational
collar and \eqref{eq:geometry:comparison-flux} on a
tube.  Closure is preserved, and the collapse
identification preserves the original real class.
The cutoff transitions remain on bounded pieces.
For a varying decomposition choose finitely many
overlapping parameter neighborhoods carrying these
constructions.  Affine combinations with a partition
of unity in the parameters remain closed on each
four-manifold, represent the same class, and have
the same prescribed middle values.  A finite
combination retains the uniform norm bounds.  This
also handles the stationary overlaps in the low-side
and small-width interpolations.

These are curvature forms, and the argument does
not trivialize a torsion line bundle on a rational
boundary.  Such a restriction admits a flat unitary
connection; using that flat reference and splicing
a real primitive changes its curvature by the
exact forms just constructed.  Fixed torsion
holonomy and determinant transition maps are
therefore retained throughout.

Four-dimensional two-form energy is invariant under
conformal change and constant metric rescaling.
The large conformal factors at the satellite and
the rescaling by its shrinking parameter thus
introduce no new energy factor.  Their fixed
truncations contribute at most $C_T$.  The long
rational middles contribute zero.  On a tube the
metric converges with the uniform bulk error
$\eta(T)$ of Lemma~\ref{lem:geometry:tube}; the fixed
vertical form has bounded product norm, so its
integral differs from the product integral by at
most $\eta(T)$ times the tube length.  Since the
total tube length is $O(\epsilon^{-1})$, the
two-stage limit in \eqref{eq:geometry:integrated-error}
proves \eqref{eq:geometry:comparison-energy}.

For scalar curvature, compact truncated pieces
have a uniform scalar bound and bounded volume.
The cylindrical annuli and exact rational
products other than the $Y_j$ products have nonnegative
scalar curvature by construction; free interpolations
occur on the bounded pieces.  Each retained $Y_j$
prefix has fixed bounded local geometry and length
$O(\epsilon^{-1/2})$, so the integral of its squared
negative scalar part is $O(\epsilon^{-1/2})
=o(\epsilon^{-1})$.  The omitted middle may be
arbitrarily long and is not included in this assertion.
On a tube, two spatial derivative
bounds control its scalar defect from the product
value four.  Its bounded defect tending to zero
in the bulk has the same estimate
\eqref{eq:geometry:integrated-error}, also after
squaring.  This proves the last assertion.
\end{proof}

In the application $\lambda_U=-1-e_j+e_{j+1}$, so
$|\lambda_U|\le2$.  The fiber has area $2\pi$, scalar
curvature four, and the circle and longitudinal
normalizations are exactly those in
\eqref{eq:geometry:long-model}.  These fixed constants
are what make the Dirac estimate compatible with
the prescribed flux.  Lemma~\ref{lem:geometry:comparison}
completes the proof of Theorem~\ref{thm:geometry:family};
the next section chooses the background connections
and then chooses finite auxiliary sizes so that all
the required tests hold simultaneously.
